% !TEX root = ../Abschlussarbeit_Jan_Händl.tex
\newglossaryentry{Express}{
  name={EXPRESS},
  description={EXPRESS ist ein Standard für eine generische Datenmodellierungssprache für Produktdaten.
      EXPRESS ist im ISO-Standard für den Austausch von Produktmodelldaten STEP (ISO 10303) formalisiert und als ISO 10303-11 standardisiert}
}
\newglossaryentry{API}{
  name={API},
  description={Eine Schnittstelle, die es zwei oder mehr Programmen ermöglicht, miteinander zu kommunizieren (Application Programming Interface)}
}

\newglossaryentry{Framework}{
  name={Framework},
  description={Ein vorgefertigtes Gerüst, das die Entwicklung von Software vereinfacht}
}

\newglossaryentry{Webapplikation}{
  name={Webanwendung},
  description={Eine Anwendung, die über einen Webbrowser zugänglich ist},
  plural={Webapplikationen}
}

\newglossaryentry{Bibliothek}{
  name={Bibliothek},
  description={Eine Bibliothek ist eine Sammlung von Funktionen oder Klassen, die eine Programmiersprache um weitere Befehle erweitert}
}

\newglossaryentry{Rest-API}{
  name={Rest-API},
  description={
      Ein Architekturstil für verteilte Systeme, der häufig für die Entwicklung von Web-APIs genutzt wird (Representational State Transfer)}
}
\newglossaryentry{Plugin}{
  name={Plugin},
  description={
      Ein Plugin ist eine Softwareerweiterung, die eine Anwendung um zusätzliche Funktionalität erweitert, ohne den Quellcode der Hauptanwendung zu ändern}
}
\newglossaryentry{acid}{
  name={ACID},
  description={Ein Akronym für die vier geforderten Eigenschaften von Datenbanktransaktionen zur Sicherstellung der Datenintegrität:
      \textbf{A}tomicity (Atomarität),
      \textbf{C}onsistency (Konsistenz),
      \textbf{I}solation (Isolation) und
      \textbf{D}urability (Dauerhaftigkeit)},
  first={ACID (Atomicity, Consistency, Isolation, Durability)}
}
\newglossaryentry{SPA}
{
  name={SPA},
  description={Eine Single-Page-Application (deutsch: Einzelseiten-Anwendung) ist eine Webapplikation, die aus einem einzigen HTML-Dokument besteht.
      Anstatt bei jeder Nutzerinteraktion die komplette Seite vom Server neu zu laden, werden Inhalte dynamisch via JavaScript (z. B. durch \textit{React}) nachgeladen und manipuliert. Dies ermöglicht eine flüssige Benutzererfahrung, die der einer nativen Desktop-Anwendung ähnelt},
  first={Single-Page-Application (SPA)},
  plural={SPAs}
}
\newglossaryentry{STEP}{
  name={STEP},
  long={Standard for the Exchange of Product Model Data},
  description={Standard for the Exchange of Product Model Data (normiert nach ISO 10303). Ein plattformunabhängiges, standardisiertes Datenformat zum Austausch von 3D-Produktdaten und Bauwerksmodellen zwischen heterogenen CAx- und BIM-Systemen.
      Im Kontext von \textit{Building Information Modeling} (BIM) bildet die physikalische Textdatei nach ISO 10303-21 (\textit{Clear Text Encoding}) das fundamentale Austauschformat für klassische ifc-Dateien}
}
\newglossaryentry{webpack}
{
  name={Webpack},
  description={Ein moderner Open-Source-Modul-Bundler für JavaScript-Anwendungen. Webpack analysiert den Modul-Graphen einer Webapplikation, löst Abhängigkeiten auf und bündelt diese in optimierte statische Assets.
      In der vorliegenden Arbeit dient Webpack als technische Basis für die Implementierung der Microfrontend-Architektur mittels \textit{Module Federation}, was das dynamische Laden von funktionalen Add-Ons zur Laufzeit ermöglicht}
}
\newglossaryentry{vcs}{
  name={Versionskontrollsystem},
  plural={Versionskontrollsysteme},
  description={Eine Software zur lückenlosen Aufzeichnung, Verwaltung und Historisierung von Änderungen an digitalen Dokumenten, primär von Quelltexten im Software-Engineering.
      Ein VCS ermöglicht die kollaborative Arbeit an einer gemeinsamen Codebasis durch Mechanismen wie Branching (Abspalten von Entwicklungszweigen) und Merging (Zusammenführen von Code)}
}
\newglossaryentry{turborepo}{
  name={Turborepo},
  description={Ein von Vercel entwickeltes, inkrementelles Hochleistungs-Build-System für JavaScript- und TypeScript-Monorepos.
      Es optimiert und beschleunigt die Orchestrierung von Tasks (wie Builds, Linting und Tests) durch intelligentes Dependency-Graph-Management, inhaltsbasiertes lokales Caching
      sowie Remote-Caching innerhalb von CI/CD-Pipelines}
}
\newglossaryentry{npm}{
  name={npm},
  description={Node Package Manager, der Standard-Paketmanager für die JavaScript-Laufzeitumgebung \texttt{Node.js} zur Verwaltung von Projektabhängigkeiten}
}
\newglossaryentry{jotai}{
  name={Jotai},
  description={Eine minimalistische und atomare State-Management-Bibliothek für React-Anwendungen.
      Der Zustand wird in kleinsten Einheiten, sogenannten \textit{Atoms}, verwaltet}
}
\newglossaryentry{JSON}{
  name={JSON},
  description={Ein leichtgewichtiges Datenformat, das zum Austausch von Daten zwischen Server und Client verwendet wird (JavaScript Object Notation)}
}
\newglossaryentry{eventemitter}{
  name={Event Emitter (Browser)},
  description={Ein Architekturmuster im Webbrowser, das die entkoppelte, ereignisgesteuerte Kommunikation
      zwischen Benutzeroberfläche und Programmcode ermöglicht.
      Es sorgt dafür, dass Webseiten dynamisch auf Benutzerinteraktionen
      (wie Klicks, Tastatureingaben oder Scrollen) sowie auf Systemereignisse (wie das Laden von Daten) reagieren können,
      ohne dass die beteiligten Programmkomponenten direkt voneinander abhängig sein müssen}
}
\newglossaryentry{Interface}{
  name={Interface},
  description={Eine formale Schnittstelle in der Softwareentwicklung,
      die als Vertrag für Klassen oder Objekte fungiert.
      Sie definiert deklarativ, welche Methoden, Eigenschaften oder Ereignisse nach außen bereitgestellt werden müssen,
      enthält jedoch keinerlei eigene Programmierlogik}
}
\newglossaryentry{middleware}{
  name={Middleware},
  description={Eine zwischengelagerte Softwarekomponente oder Softwareschicht,
      die als Bindeglied zwischen verschiedenen Anwendungen, Systemen oder Prozessen agiert}}
\newglossaryentry{local_storage}{
  name={Local Storage},
  description={Ein permanenter Web-Speicher-Mechanismus des Browsers, der es Webanwendungen ermöglicht,
      Daten in Form von Key-Value-Paaren lokal auf dem Endgerät des Nutzers zu speichern.
      Im Gegensatz zu Cookies werden die Daten nicht bei jeder HTTP-Anfrage an den Server übertragen und im Gegensatz
      zum Session-Storage bleiben die Daten auch nach dem Schließen des Browser-Tabs oder des Fensters dauerhaft erhalten,
      bis sie programmatisch oder vom Nutzer gelöscht werden}
}

\newglossaryentry{docker}{
  name={Docker},
  description={Eine Open-Source-Plattform zur Containerisierung von Anwendungen. Docker ermöglicht es, Anwendungen inklusive ihrer Abhängigkeiten, Bibliotheken und Konfigurationsdateien in isolierte, portable Container zu verpacken. Im Gegensatz zu traditionellen virtuellen Maschinen teilen sich Docker-Container den Kernel des Host-Betriebssystems, wodurch sie ressourcenschonender, kompakter und schneller im Startverhalten sind. Dies stellt eine konsistente Laufzeitumgebung über den gesamten Entwicklungs- und Bereitstellungszyklus hinweg sicher}
}

\newglossaryentry{docker_compose}{
  name={Docker Compose},
  description={Ein Werkzeug zur Definition und Ausführung von Multi-Container-Anwendungen in Docker. Über eine Konfigurationsdatei (typischerweise \texttt{docker-compose.yml}) im YAML-Format können alle benötigten Dienste, Netzwerke und Daten-Volumes einer Anwendung deklariert und anschließend mit einem einzigen Befehl koordiniert gestartet, gestoppt und verwaltet werden}
}
